\chapter{کتابخانهٔ استاندارد}
\label{ch:stdlib}

هیچ برنامه‌نویسی همه‌چیز را از صفر نمی‌سازد. سلام با یک \textbf{کتابخانهٔ
استاندارد} غنی همراه است: مجموعه‌ای از ماژول‌های آماده برای ریاضی، متن،
ساختارهای داده، زمان، تصادف و بسیار چیزهای دیگر. در این فصل پرکاربردترین‌ها را
می‌بینیم.

\section{ریاضی}

ماژولِ \code{ریاضی} (math) توابعِ ریاضیِ رایج را در اختیار می‌گذارد. این مثالِ
واقعی چند تای آن‌ها را نشان می‌دهد:

\begin{salam}
واردکردن "math"
تابع اصلی:
    چاپ("جذر ۱۴۴:", ریاضی.جذر(144.0))
    چاپ("بیشینه:", ریاضی.بیشینه(3.0, 9.0))
    چاپ("عدد پی:", ریاضی.پی)
پایان
\end{salam}

مهم‌ترین توابع و ثابت‌های ماژولِ ریاضی:

\begin{center}
\begin{tabular}{l l l}
\toprule
\textbf{فارسی} & \textbf{انگلیسی} & \textbf{کار} \\
\midrule
\code{ریاضی.جذر(x)} & \code{Sqrt} & جذر \\
\code{ریاضی.توان(x, y)} & \code{Pow} & توان \\
\code{ریاضی.قدرمطلق(x)} & \code{Abs} & قدرِمطلق \\
\code{ریاضی.کف(x)} & \code{Floor} & گردکردن به پایین \\
\code{ریاضی.سقف(x)} & \code{Ceil} & گردکردن به بالا \\
\code{ریاضی.سینوس(x)} & \code{Sin} & سینوس \\
\code{ریاضی.کسینوس(x)} & \code{Cos} & کسینوس \\
\code{ریاضی.لگاریتم(x)} & \code{Log} & لگاریتم \\
\code{ریاضی.بیشینه(a, b)} & \code{Max} & بزرگ‌تر \\
\code{ریاضی.پی} & \code{PI} & عددِ $\pi$ \\
\bottomrule
\end{tabular}
\end{center}

\begin{warn}
بیشترِ توابعِ ریاضی روی اعدادِ \emph{اعشاری} کار می‌کنند. به همین دلیل
\code{ریاضی.جذر(144.0)} نوشتیم، نه \code{ریاضی.جذر(144)}. اگر عددِ صحیح دارید،
نخست آن را با \code{بعنوان اعشاری۶۴} تبدیل کنید.
\end{warn}

\section{رشته}

ماژولِ \code{رشته} (str) را در فصلِ ۸ معرفی کردیم. علاوه بر متدهای پایه، توابعِ
پیشرفته‌تری نیز دارد. این مثالِ واقعی چند تای آن‌ها را نشان می‌دهد:

\begin{salam}
تابع اصلی:
    ثابت متن رشته = "salam,donya"
    چاپ "طول-متن", متن.طول()
    چاپ "بخش", متن.زیررشته(0, 5)
    چاپ "یافتن", متن.بیاب("donya")
    چاپ "پیوست", متن.پیوست("!")
    متغیر بخش‌ها وکتور<رشته> = متن.بشکاف(",")
    چاپ "تعداد-بخش", بخش‌ها.طول()
    چاپ "عدد", "۴۲ ".پیراست()
    چاپ "به‌صحیح", "42".به_صحیح() + 1
پایان
\end{salam}

توابعِ تازه‌ای که اینجا می‌بینیم:
\begin{itemize}
  \item \code{.بیاب("...")} (\code{Find}): جای نخستین یافتنِ یک زیررشته را
        برمی‌گرداند.
  \item \code{.بشکاف(",")} (\code{Split}): رشته را بر اساسِ یک جداکننده به یک
        \emph{بردار} از تکه‌ها می‌شکند.
  \item \code{.پیراست()} (\code{Trim}): فاصله‌های آغاز و پایان را می‌زداید.
  \item \code{.به\_صحیح()} (\code{ToInt}): یک رشتهٔ عددی را به عددِ صحیح تبدیل
        می‌کند.
\end{itemize}

\begin{tip}
\code{.بشکاف()} یکی از پرکاربردترین توابع در پردازشِ داده است. هر وقت داده‌ای
دارید که با کاما، فاصله یا نشانهٔ دیگری جدا شده (مثلِ یک خطِ فایلِ CSV)، با
\code{بشکاف} آن را به اجزایش تقسیم کنید.
\end{tip}

\section{نگاشت: واژه‌نامهٔ کلید-مقدار}

\textbf{نگاشت} (HashMap) ساختاری است که هر \emph{کلید} را به یک \emph{مقدار}
وصل می‌کند درست مانندِ یک واژه‌نامه که هر واژه را به معنایش وصل می‌کند. این
مثالِ واقعی سن‌های چند نفر را نگه می‌دارد:

\begin{salam}
تابع اصلی:
    متغیر سن‌ها نگاشت<رشته, صحیح۳۲> = نگاشت {}
    سن‌ها.درج("سارا", 30)
    سن‌ها.درج("علی", 25)
    سن‌ها.درج("رضا", 40)
    چاپ("سن سارا:", سن‌ها.بگیر("سارا"))
    چاپ("علی موجود است؟", سن‌ها.دارد("علی"))
    چاپ("حسن موجود است؟", سن‌ها.دارد("حسن"))
    چاپ("تعداد:", سن‌ها.اندازه())
پایان
\end{salam}

نوعِ \code{نگاشت<رشته, صحیح۳۲>} یعنی «نگاشتی که کلیدهایش رشته و مقدارهایش عددِ
صحیح‌اند». متدهای اصلی:

\begin{center}
\begin{tabular}{l l l}
\toprule
\textbf{فارسی} & \textbf{انگلیسی} & \textbf{کار} \\
\midrule
\code{.درج(کلید, مقدار)} & \code{put} & افزودن یا به‌روزرسانی \\
\code{.بگیر(کلید)} & \code{get} & گرفتنِ مقدارِ یک کلید \\
\code{.دارد(کلید)} & \code{has} & آیا کلید هست؟ \\
\code{.حذف(کلید)} & \code{remove} & حذفِ یک کلید \\
\code{.اندازه()} & \code{size} & تعدادِ عناصر \\
\code{.پیمایش()} & \code{iter} & ساختنِ پیمایشگر \\
\bottomrule
\end{tabular}
\end{center}

\subsection{پیمایشِ یک نگاشت}

برای گذشتن از همهٔ عناصرِ یک نگاشت، یک \textbf{پیمایشگر} می‌سازیم. این مثالِ
واقعی، مقدارهای یک نگاشت را جمع می‌زند:

\begin{salam}
تابع اصلی:
    متغیر امتیاز نگاشت<رشته, صحیح۳۲> = نگاشت {}
    امتیاز.درج("الف", 3)
    امتیاز.درج("ب", 7)
    متغیر مجموع صحیح۳۲ = 0
    متغیر پیمایه پیمایشگر_نگاشت<رشته, صحیح۳۲> = امتیاز.پیمایش()
    تاوقتی پیمایه.دارد_بعدی():
        مجموع = مجموع + پیمایه.مقدار()
        پیمایه.بعدی()
    پایان
    چاپ "جمع", مجموع
پایان
\end{salam}

پیمایشگر سه متدِ کلیدی دارد: \code{دارد\_بعدی()} (آیا عنصرِ دیگری مانده؟)،
\code{مقدار()} (مقدارِ عنصرِ جاری) و \code{بعدی()} (رفتن به عنصرِ بعد). این الگو
--- «تا وقتی عنصری مانده، پردازش کن و جلو برو» شیوهٔ استانداردِ گذشتن از
مجموعه‌هاست.

\begin{tip}
نگاشت برای جست‌وجوی سریع عالی است. اگر می‌خواهید بدانید «سنِ سارا چند است؟»،
نگاشت پاسخ را تقریباً بی‌درنگ می‌دهد بی‌آنکه مجبور باشید مثلِ یک بردار،
همه‌چیز را یکی‌یکی بگردید. هر جا نیاز به نگاشتنِ «نام به داده» داشتید (شمارهٔ
دانشجو به نمره، نامِ کالا به قیمت)، نگاشت ابزارِ درست است.
\end{tip}

\section{دیگر ماژول‌های کاربردی}

کتابخانهٔ استانداردِ سلام بسیار گسترده است. در ادامه چند ماژولِ پرکاربردِ دیگر
را به‌اختصار معرفی می‌کنیم؛ شیوهٔ کارشان همان است: \code{واردکردن}، سپس دسترسی
با نامِ ماژول و نقطه.

\begin{center}
\begin{tabular}{l l}
\toprule
\textbf{ماژول} & \textbf{کاربرد} \\
\midrule
\code{io} & ورودی و خروجی؛ خواندنِ ورودیِ کاربر، چاپ، کار با فایل \\
\code{time} & زمان و تاریخ؛ \code{Now}، سال، ماه، روز، \code{Sleep} \\
\code{rand} & تولیدِ عددهای تصادفی؛ \code{Seed}، \code{Int}، \code{Float} \\
\code{fmt} & قالب‌بندیِ متن؛ ساختنِ رشته از عدد و مقدارهای دیگر \\
\code{os} & سیستم‌عامل؛ آرگومان‌ها، متغیرهای محیطی، فایل‌ها \\
\code{json} & خواندن و نوشتنِ داده‌های JSON \\
\code{testing} & نوشتنِ آزمون؛ \code{Assert} و گزارشِ نتیجه \\
\code{sync} & هم‌روندی؛ نخ‌ها و قفل‌ها (فصلِ ۱۷) \\
\code{http} & درخواست‌های وب؛ \code{Get}، \code{Post} \\
\bottomrule
\end{tabular}
\end{center}

ماژول‌های \code{os}، \code{io}، \code{time} و \code{http} را در فصل‌های آینده
(۱۶ و ۱۷) بیشتر می‌بینیم؛ \code{testing} را در فصلِ ۲۰.

\subsection{نمونه: عددِ تصادفی}

\begin{salam}
واردکردن "rand"
تابع اصلی:
    rand.Seed(7)
    چاپ "یک عددِ تصادفی:", rand.Int()
پایان
\end{salam}

نخست با \code{Seed} «بذرِ» مولدِ تصادف را می‌کاریم (تا دنبالهٔ اعداد قابلِ
بازتولید باشد)، سپس با \code{Int} یک عددِ تصادفی می‌گیریم.

\begin{deep}[نگاهی به درون: تصادفِ واقعی یا ساختگی؟]
رایانه نمی‌تواند عددِ «واقعاً» تصادفی بسازد؛ آنچه می‌سازد دنباله‌ای است که
\emph{تصادفی به‌نظر می‌رسد} اما در واقع از یک فرمولِ معین و یک «بذرِ» آغازین
زاده می‌شود. به همین دلیل اگر بذرِ یکسان بدهید، همان دنباله را می‌گیرید که
برای آزمون و اشکال‌زدایی بسیار مفید است. اگر تصادفِ غیرقابلِ‌پیش‌بینی
می‌خواهید، بذر را از زمانِ جاری بگیرید.
\end{deep}

\section{خلاصهٔ فصل}

\begin{itemize}
  \item ماژولِ \code{ریاضی} توابعِ ریاضی و ثابت‌هایی مانندِ $\pi$ دارد.
  \item ماژولِ \code{رشته} توابعِ پیشرفتهٔ متن دارد: \code{بشکاف}،
        \code{بیاب}، \code{پیراست}، \code{به\_صحیح}.
  \item \code{نگاشت} هر کلید را به یک مقدار وصل می‌کند و جست‌وجو را سریع
        می‌سازد.
  \item پیمایشگر با \code{دارد\_بعدی}، \code{مقدار} و \code{بعدی} از مجموعه‌ها
        می‌گذرد.
  \item ماژول‌های فراوانِ دیگری برای زمان، تصادف، فایل، شبکه و آزمون هست.
\end{itemize}

\section{تمرین‌ها}

\exercise با ماژولِ \code{ریاضی}، طولِ وترِ مثلثی با اضلاعِ ۳ و ۴ را حساب کنید
(قضیهٔ فیثاغورس).

\exercise یک نگاشت بسازید که نامِ چند میوه را به قیمتشان وصل کند، سپس قیمتِ یک
میوه را بپرسید.

\exercise با پیمایشِ یک نگاشتِ قیمت‌ها، گران‌ترین کالا را پیدا کنید.

\exercise رشتهٔ \code{"۱۰،۲۰،۳۰،۴۰"} را با \code{بشکاف} به اعداد تبدیل کنید و
مجموعشان را بیابید.

\exercise با ماژولِ \code{rand}، یک «تاسِ» شش‌وجهی بسازید که عددی میانِ ۱ تا ۶
بدهد (راهنمایی: از باقی‌ماندهٔ تقسیم بر ۶ کمک بگیرید).
